\begingroup\footnotesize
\begin{table}
\centering
\caption{\label{tab:fishery_stats}Commercial and Recreational fisheries in Gulf States are an important source of livelihoods and food.}
\centering
\begin{tabular}[t]{lrrrrrrrrrr}
\toprule
\multicolumn{1}{c}{ } & \multicolumn{2}{c}{Employment} & \multicolumn{2}{c}{Production} & \multicolumn{2}{c}{Income} & \multicolumn{2}{c}{Sales} & \multicolumn{2}{c}{Value Added} \\
\cmidrule(l{3pt}r{3pt}){2-3} \cmidrule(l{3pt}r{3pt}){4-5} \cmidrule(l{3pt}r{3pt}){6-7} \cmidrule(l{3pt}r{3pt}){8-9} \cmidrule(l{3pt}r{3pt}){10-11}
State & Com. & Rec. & Com. & Rec. & Com. & Rec. & Com. & Rec. & Com. & Rec.\\
\midrule
Florida & 121,710 & 28,546 & 49,902 & 60,556 & 4,579 & 1,350 & 24,596 & 3,936 & 8,207 & 2,401\\
Texas & 41,171 & 5,605 & 27,889 &  & 1,259 & 261 & 5,584 & 793 & 2,049 & 450\\
Louisiana & 32,514 & 4,552 & 413,855 &  & 623 & 171 & 1,638 & 590 & 838 & 294\\
Alabama & 6,971 & 5,506 & 13,872 & 8,685 & 140 & 217 & 443 & 745 & 196 & 399\\
Mississippi & 6,954 & 324 & 141,018 & 4,917 & 127 & 13 & 329 & 72 & 165 & 30\\
\bottomrule
\multicolumn{11}{l}{\rule{0pt}{1em}Data come from citep{noaa_feus_2022}. Each economic metric is split by fishing sector: Com. = commercial fisheries; Rec. = recreational fisheries. Units: Employment is full- and part-time jobs; Production is in metric tons; Income, Sales and Value Added are in millions of U.S. dollars (M USD). Empty cells indicate that the metric is not separately reported for that state-sector combination.}\\
\end{tabular}
\end{table}
\endgroup
